%! Describing the Distribution of a Quantitative Variable — Unit 1, Lesson 1.4 — AP Practice (Answer Key)
% GRADUAL RELEASE: AP-format practice, extra credit, scored 1 to 5 (the score is added as extra-credit points). The graded practice is the `homework` component that
% follows it at the back of the packet.
% MC 2 carries the target misconception as distractor (B) ("skewed left because the pile is on
% the left"); MC 3 carries the EK UNC-1.H.7 over-claim. Free-response (a) spirals to Lesson 1.1
% (classify the variable first) and (d) to Lesson 1.3 (check the interval width before you
% believe the shape).
% GENERATED FROM ap_practice/main.tex. Each \writelines{n} in the blank is answered by an
% \ansline{} terse enough not to wrap past n lines — that keeps the page counts equal.
% NO teachernote here — scoring guidance lives in the lesson plan.
%
% THE DISPLAY: 40 commute times, 10-minute intervals from 0: 5 12 10 6 4 2 0 0 1 = 40.
% Right-skewed, with one employee stranded in the 80-90 interval behind a two-interval gap.
% Regrouped at 30 minutes: 27 / 12 / 1 — the gap vanishes; the long commute is alone in 60-90.
\documentclass[12pt]{article}
\usepackage{apstats-article}
\usepackage{apstats-key}   % pulls in -boxes; do NOT also load -boxes
\usepackage{grading}
\gradingassignment{apstats-u01-l04-ap}
\gradingcourse{advanced_statistics}

\begin{document}

\pageheader{Unit 1, Lesson 1.4}{AP Practice --- Answer Key}
{\small\itshape Write each final answer in its box. Only what is inside a box is graded.}\par
% NAMESTRIP: no \namedateperiod here — mirrors the blank.

\vspace{0.08in}

% Authored BYTE-IDENTICALLY in ap_practice/ and ap_practice_key/ — this box is what keeps the
% blank and the key the same number of pages.
\boxguard[8]
\begin{remindbox}
\small
This page is \textbf{extra credit, scored 1 to 5} --- your graded homework is the last section
of this packet. These are AP-format reps: work them for the practice, and bring what you get
stuck on to office hours or study hall.
\end{remindbox}

\vspace{0.12in}

\boxguard[24]
\begin{scenariobox}[The Commute Record]{navy}
\small
A company recorded the one-way commute time, in minutes, of each of its 40 employees. The
histogram below groups those 40 commute times into intervals 10 minutes wide.

\vspace{6pt}
\begin{center}
\begin{tikzpicture}[scale=0.80]
  \foreach \i/\c in {0/5, 1/12, 2/10, 3/6, 4/4, 5/2, 6/0, 7/0, 8/1} {
    \fill[navy] ({\i*0.6},0) rectangle ({(\i+1)*0.6},{\c*0.18});
    \draw[white, very thin] ({\i*0.6},0) rectangle ({(\i+1)*0.6},{\c*0.18});}
  \draw[->, thick] (0,0) -- (0,2.7);
  \draw[->, thick] (0,0) -- (5.9,0);
  \foreach \y/\lab in {0/0, 0.9/5, 1.8/10} {
    \draw (-0.07,\y) -- (0.07,\y); \node[left, font=\tiny, xshift=-1pt] at (-0.07,\y) {\lab};}
  \foreach \x/\lab in {0/0, 1.2/20, 2.4/40, 3.6/60, 4.8/80} {
    \draw (\x,-0.07) -- (\x,0.07); \node[below, font=\tiny] at (\x,-0.1) {\lab};}
  \node[font=\tiny] at (2.7,-0.85) {One-way commute time (minutes)};
  \node[font=\tiny, rotate=90] at (-0.62,1.2) {Employees};
\end{tikzpicture}
\end{center}
\end{scenariobox}

\vspace{0.10in}

\boxguard[16]
\begin{headlinebox}{goldbg}
\small
\textbf{Section I --- Multiple Choice.} Write the letter of the single best answer in its Letter box. Every answer choice is
about the context, not just the arithmetic --- read them all before choosing.
\end{headlinebox}

\vspace{0.10in}

\begin{enumerate}[leftmargin=1.9em, itemsep=0.11in]

  \item How many of these 40 employees have a one-way commute of 30 minutes or more?
        \begin{enumerate}[label=\textbf{(\Alph*)}, leftmargin=2.2em, itemsep=2pt]
          \item 6
          \item 10
          \item \ans{13}
          \item 17
          \item 27
                \hfill\gradedright{\textit{Letter:}~\gradedblank[id=ap1, type=choice, points=1, answer=C]{0.45in}{0.24in}}
        \end{enumerate}

  \item Which statement best describes the shape of this distribution?
        \begin{enumerate}[label=\textbf{(\Alph*)}, leftmargin=2.2em, itemsep=2pt]
          \item Approximately symmetric and unimodal.
          \item Skewed to the left, because most of the commutes lie on the left side of the
                graph.
          \item \ans{Skewed to the right, because the longer commutes trail off to the right.}
          \item Approximately uniform, because every interval contains at least a few employees.
          \item Bimodal, because two bars are taller than the rest.
                \hfill\gradedright{\textit{Letter:}~\gradedblank[id=ap2, type=choice, points=1, answer=C]{0.45in}{0.24in}}
        \end{enumerate}

  \boxguard[10]
  \item Which conclusion is \textbf{supported} by this display?
        \begin{enumerate}[label=\textbf{(\Alph*)}, leftmargin=2.2em, itemsep=2pt]
          \item Employees at this company commute longer than employees at other companies.
          \item Commute times for all workers in this city are skewed to the right.
          \item \textcolor{keyred}{Most of these 40 employees commute under 30 minutes; one commutes over 80.}  % not bold: bold wraps and breaks the mirror
          \item The employee with the longest commute should be removed before the distribution
                is described.
          \item Half of these employees commute more than 40 minutes.
                \hfill\gradedright{\textit{Letter:}~\gradedblank[id=ap3, type=choice, points=1, answer=C]{0.45in}{0.24in}}
        \end{enumerate}

  \item Before concluding that this distribution is skewed to the right, a careful reader should
        first check
        \begin{enumerate}[label=\textbf{(\Alph*)}, leftmargin=2.2em, itemsep=2pt]
          \item whether commute time is a categorical variable.
          \item \ans{the interval width, because a different width can change the shape.}
          \item whether the bars have gaps between them, as in a bar graph.
          \item whether any employee was counted twice.
          \item nothing --- a histogram reports the shape of the data as a fact.
                \hfill\gradedright{\textit{Letter:}~\gradedblank[id=ap4, type=choice, points=1, answer=B]{0.45in}{0.24in}}
        \end{enumerate}

\end{enumerate}

\vspace{0.10in}

\boxguard[16]
\begin{headlinebox}{goldbg}
\small
\textbf{Section II --- Free Response.} Show your reasoning. Every answer must be written
\emph{in the context of the problem} --- that is what earns credit on the AP exam.
\end{headlinebox}

\vspace{0.10in}

\begin{enumerate}[leftmargin=1.9em, itemsep=0.12in, resume]

  \item Use the commute record above.

        \boxguard[17]
        \vspace{4pt}
        \textbf{(a)} Classify commute time as categorical or quantitative, and then as discrete
        or continuous. Justify from the variable itself, not from the display.\\[4pt]
        \gradedhere{\gradedblank[id=ap5a, type=rubric, points=2]{\linewidth+6pt}{\gradedlinesheight{2}}}\ansline{Quantitative --- a commute time is a measured amount, not a category. It is}\\[1pt]
        \ansline{continuous: between 20 and 21 minutes, 20.4 minutes is perfectly possible.}\\
        \gradedhide{\writelines{7}}

        \boxguard[17]
        \vspace{4pt}
        \textbf{(b)} Describe the shape of this distribution and identify any unusual features,
        in context.\\[4pt]
        \gradedhere{\gradedblank[id=ap5b, type=rubric, points=2]{\linewidth+6pt}{\gradedlinesheight{2}}}\ansline{Skewed right: most of the 40 employees commute under 30 minutes, with a tail of}\\[1pt]
        \ansline{longer ones. One is an outlier at 80--90 minutes, behind a 20-minute empty stretch.}\\
        \gradedhide{\writelines{7}}

        \boxguard[17]
        \vspace{4pt}
        \textbf{(c)} A manager writes: \emph{``The average commute here is about 25 minutes, so
        our employees are fine.''} Explain what that single number leaves out, and state what
        these data can and cannot tell you about employees at other companies.\\[4pt]
        \gradedhere{\gradedblank[id=ap5c, type=rubric, points=2]{\linewidth+6pt}{\gradedlinesheight{3}}}\ansline{It leaves out shape, spread, and the unusual value: these commutes run from under}\\[1pt]
        \ansline{10 minutes to over 80, and one employee travels more than 80 minutes each way.}\\[1pt]
        \ansline{These data describe only these 40 employees, not workers at any other company.}\\
        \gradedhide{\writelines{6}}

        \boxguard[17]
        \vspace{4pt}
        \textbf{(d)} The same 40 commute times are redrawn using intervals 30 minutes wide.
        Explain what happens to the employee who commutes more than 80 minutes, and why the
        interval width has to be checked before the shape is believed.\\[4pt]
        \gradedhere{\gradedblank[id=ap5d, type=rubric, points=2]{\linewidth+6pt}{\gradedlinesheight{2}}}\ansline{The bars become 27, 12, 1. That employee is alone in a 60--90 bar that now sits}\\[1pt]
        \ansline{right beside the 30--60 bar: the gap vanished, but the data did not change.}\\
        \gradedhide{\writelines{7}}

\end{enumerate}


\end{document}
